% Section: Budget Pacing Under Cost Drift
% Include from paper/sections/results.tex

\subsection{Budget Pacing Under Cost Drift (Figure~\ref{fig:exp03_adaptation})}
\label{sec:budget_drift}

\paragraph{Motivation.}
Experiments 01 and 02 studied budget pacing and non-stationarity in
isolation.  In production, both occur simultaneously: a router
operates under a cost budget when a model provider drops its pricing.
We ask: \emph{does the BudgetPacer automatically exploit a mid-stream
price drop to improve quality while maintaining budget compliance,
and how does it compare with static cost-control baselines?}

\paragraph{Setup.}
The $K{=}3$ portfolio follows a \emph{train-then-evaluate} pipeline
(matching Experiment~01).
The router first online-learns on the validation split
($n{=}1{,}785$) under normal pricing without recording evaluation
metrics; priors are trained on the disjoint training split
($n_{\mathrm{eff}}{=}\bdNeff$, $\alpha{=}0.1$;
Section~\ref{sec:hparam_sweep}).
Evaluation uses the held-out test split ($n{=}1{,}824$),
partitioned into three phases of $608$ steps each
(phase boundaries at steps $608$ and $1{,}216$):
Phase~1 uses normal pricing (Gemini normalized cost
$\tilde{c}{=}0.67$); Phase~2 applies a Gemini price
drop to \$0.10/\$0.10 per million tokens ($\tilde{c} \approx 0.0$);
Phase~3 restores the original Gemini pricing
($\tilde{c}{=}0.67$).
Phase~3 deliberately reuses the same 608 prompts as Phase~1 so that
the P1-vs-P3 comparison is a controlled within-subject design: any
difference between the two phases is attributable solely to the
router's internal state (accumulated Phase~2 learning history), not
prompt-sampling variability.
With geometric forgetting $\gamma{=}0.995$, Phase~1 observations are
down-weighted by $\gamma^{608}\approx 0.047$ by the start of Phase~3,
so the learner retains negligible memory of individual Phase~1
evaluations.
The Fixed Policy (no online learning) serves as a null check: its P1
and P3 metrics are identical, confirming that any P1-vs-P3 gap in
other conditions is algorithmic rather than due to prompt sampling.

Three budget targets are chosen to span qualitatively distinct
operating regimes of the primal-dual controller:
\begin{itemize}[nosep,leftmargin=*]
  \item \textbf{Tight} ($B{=}\$3.0{\times}10^{-4}$/req):
    The budget constraint is binding in Phase~1, producing a
    mixed Llama/Mistral policy with high $\lambda_t$.
    The target sits just above the post-drop cost floor,
    so the constraint relaxes in Phase~2.
    This regime tests whether the pacer smoothly transitions
    from a binding to a slack constraint after the price change,
    and whether the constraint re-binds when pricing is restored
    in Phase~3.
  \item \textbf{Moderate} ($B{=}\$6.6{\times}10^{-4}$/req):
    The constraint is actively binding but not saturating.
    The router must continuously balance quality and cost,
    producing a mixed Llama/Mistral policy.  This is the
    regime where the closed-loop nature of the BudgetPacer
    matters most, because small routing changes cause immediate
    cost feedback that the dual update must absorb.
    Phase~3 tests whether the pacer recovers budget compliance
    when pricing reverts.
  \item \textbf{Loose} ($B{=}\$1.9{\times}10^{-3}$/req):
    The budget barely constrains routing; $\lambda_t$ is low
    and the router operates close to its unconstrained optimum.
    This regime tests whether the pacer gracefully ``gets out
    of the way'' when the budget is not a binding constraint,
    and whether it re-engages when Phase~3 restores the
    original cost structure.
\end{itemize}

\noindent
Together, these three levels cover the full spectrum from
budget-dominated to quality-dominated routing, allowing us to
evaluate adaptation behavior across the entire operating envelope
rather than at a single hand-picked budget.

\paragraph{Why this matters.}
The scenario above is not hypothetical: Gemini's pricing dropped
$12\times$ between 2024 and 2025, and similar price adjustments
occur quarterly across major providers.
A router that cannot exploit a mid-stream price drop
leaves money on the table---or, worse, continues routing to
expensive models while a cheaper, equally capable alternative
exists.
The question is whether the router's cost-control mechanism can
adapt \emph{automatically}, without operator intervention, and
whether it does so while maintaining budget compliance.

\paragraph{Why these baselines.}
We select four conditions that isolate the contribution of each
cost-control mechanism.  Each adds exactly one capability:

\noindent
Per target, four conditions represent increasing levels of
routing sophistication, each using warmup priors:
\begin{enumerate}[nosep,leftmargin=*]
  \item \textbf{Fixed Policy (offline):}
    Frozen warmup priors with a matched static cost penalty;
    no online learning.  The industry standard: deploy an
    offline-trained router with a hand-tuned cost weight.
    This establishes the \emph{cost of not adapting} under
    pricing changes.
  \item \textbf{Naive Bandit ($\gamma{=}1.0$):}
    LinUCB with infinite memory and the same matched static
    cost penalty.  Online learning, but no principled budget
    mechanism and no forgetting.
    This isolates the value of \emph{online quality learning}
    when cost control remains open-loop.
  \item \textbf{Recalibrated Bandit ($\gamma{=}1.0$):}
    Identical to the Naive Bandit in Phase~1.  At each phase
    boundary (steps $608$ and $1{,}216$), the static cost penalty
    is re-tuned offline: a grid search over $\lambda_s$ candidates
    on the validation split with the new phase's pricing selects
    the value that minimises
    $|\overline{c} - B|$.  This baseline directly isolates
    \emph{continuous online tracking} (ParetoBandit) from
    \emph{stepwise offline recalibration}, and is the strongest
    static-penalty baseline in our comparison.
    It answers the question: ``if we re-tuned the cost weight
    whenever prices change, would that be enough?''
  \item \textbf{ParetoBandit (Pacer + $\gamma{=}0.995$):}
    Warmup priors with jointly-tuned forgetting and the
    primal-dual BudgetPacer targeting the dollar budget directly.
    This is the only condition with \emph{closed-loop} cost
    control that continuously adapts to realized spend.
\end{enumerate}

\noindent
An \textbf{unconstrained ParetoBandit} baseline (same $\gamma{=}0.995$,
same priors, no budget pacer, $\lambda_s{=}0$) is included in
every panel of Figure~\ref{fig:exp03_adaptation}.
Because it shares the forgetting factor with ParetoBandit, any gap
between a constrained curve and this baseline is attributable solely
to budget enforcement---not to differences in the learning dynamics.
This clean ablation answers two questions:
(i)~\emph{how much quality does the pacer itself cost?}---the gap
quantifies the price of budget compliance in isolation; and
(ii)~\emph{does the pacer converge toward unconstrained behavior
after the price drop?}---when Gemini becomes nearly free, a
well-adapted pacer should approach the unconstrained routing mix,
and this trace provides the reference to verify convergence.

\paragraph{Fixed Policy cost awareness.}
The ``Fixed Policy'' freezes the \emph{reward model} (no online
learning) but retains live cost awareness: the router recomputes
normalized costs from the model registry at each step.
When the Gemini price drop is applied, all conditions---including
the Fixed Policy---observe the new pricing, mirroring the real-world
scenario in which price changes are publicly announced.
The Fixed Policy can therefore shift routing in response to cost
signals (e.g., from Llama to Gemini after the drop) but
\emph{cannot} update its quality estimates for any arm.

\paragraph{Static cost-penalty calibration.}
The matched static cost penalties ($\lambda_s \in \{0.10, 0.30,
0.40\}$) were calibrated against the BudgetPacer's Phase~1 spend
profile.
Because the Fixed Policy and Naive Bandit route differently from
the pacer, the same penalty produces different actual spend;
this is a fundamental limitation of open-loop cost control and is
the central result of this experiment rather than a confound.
A grid search over $\lambda_s$ per condition would require running
each condition first---precisely the closed-loop calibration that
the BudgetPacer provides automatically.

% =====================================================================
% BUDGET COMPLIANCE (THE KEY DIFFERENTIATOR)
% =====================================================================

\paragraph{Budget compliance is the key differentiator.}
The defining advantage of the BudgetPacer is \emph{precise budget
targeting}:

\begin{center}
\small
\begin{tabular}{@{}llrr@{}}
\toprule
\textbf{Budget} & \textbf{Condition} & \textbf{Phase~1 Cost} & \textbf{Ratio} \\
\midrule
\multirow{4}{*}{Tight ($\$3.0 \times 10^{-4}$)}
  & Fixed Policy     & \bdFixedTightPhaseOneCost & $\bdFixedTightPhaseOneRatio$ \\
  & Naive Bandit     & \bdNaiveTightPhaseOneCost & $\bdNaiveTightPhaseOneRatio$ \\
  & Recalibrated     & \bdRecalTightPhaseOneCost & $\bdRecalTightPhaseOneRatio$ \\
  & \textbf{ParetoBandit} & \textbf{\bdParetoBanditTightPhaseOneCost} & $\mathbf{\bdParetoBanditTightPhaseOneRatio}$ \\[3pt]
\multirow{4}{*}{Moderate ($\$6.6 \times 10^{-4}$)}
  & Fixed Policy     & \bdFixedModPhaseOneCost & $\bdFixedModPhaseOneRatio$ \\
  & Naive Bandit     & \bdNaiveModPhaseOneCost & $\bdNaiveModPhaseOneRatio$ \\
  & Recalibrated     & \bdRecalModPhaseOneCost & $\bdRecalModPhaseOneRatio$ \\
  & \textbf{ParetoBandit} & \textbf{\bdParetoBanditModPhaseOneCost} & $\mathbf{\bdParetoBanditModPhaseOneRatio}$ \\[3pt]
\multirow{4}{*}{Loose ($\$1.9 \times 10^{-3}$)}
  & Fixed Policy     & \bdFixedLoosePhaseOneCost & $\bdFixedLoosePhaseOneRatio$ \\
  & Naive Bandit     & \bdNaiveLoosePhaseOneCost & $\bdNaiveLoosePhaseOneRatio$ \\
  & Recalibrated     & \bdRecalLoosePhaseOneCost & $\bdRecalLoosePhaseOneRatio$ \\
  & \textbf{ParetoBandit} & \textbf{\bdParetoBanditLoosePhaseOneCost} & $\mathbf{\bdParetoBanditLoosePhaseOneRatio}$ \\
\bottomrule
\end{tabular}
\end{center}

\noindent
The \textbf{Ratio} column shows Phase~1 cost as a multiple of the
budget target.  The Recalibrated Bandit is identical to the Naive
Bandit in Phase~1 (recalibration occurs at the Phase~2 boundary),
so their Phase~1 costs match by construction.
Three failure modes emerge:

\begin{itemize}[nosep,leftmargin=*]
  \item \textbf{Fixed Policy under-spends} ($0.46$--$0.51\times$).
    The static cost penalty, tuned to produce similar Phase~1 spend
    for the pacer conditions, is too aggressive for a frozen policy
    that cannot refine its reward estimates.
    At tight budget, it routes ${\sim}72\%$ to Llama-8B---the cheapest arm
    but also the lowest-quality---wasting roughly half the quality
    headroom that the budget allows.

  \item \textbf{Naive Bandit is inconsistent across budgets.}
    After pre-training on the validation split, the Naive Bandit
    overshoots at tight ($1.52\times$) and moderate ($1.16\times$)
    because the static cost penalty provides no feedback mechanism to
    correct realised over-spend.
    At loose budget, the same static penalty ($\lambda_s{=}0.10$)
    cannot contain the shift toward expensive arms: the Naive Bandit
    spends \$5.1e-3/req ($2.74\times$ target).
    A different penalty per budget tier would require the very
    closed-loop calibration that the BudgetPacer provides automatically.

  \item \textbf{Recalibrated Bandit inherits Phase~1 weakness.}
    Because recalibration only takes effect at the Phase~2 boundary,
    the Recalibrated Bandit offers no Phase~1 improvement.
    Any production scenario where budget compliance matters from
    deployment onward---rather than only after a price change---cannot
    be addressed by periodic offline recalibration.
\end{itemize}

\noindent
\textbf{ParetoBandit is the only condition that reliably hits the
budget target} ($0.97$--$1.00\times$).  The BudgetPacer's
primal-dual mechanism adjusts $\lambda_t$ continuously based on
observed costs, providing a closed-loop budget controller that
static penalties cannot replicate.

% =====================================================================
% ADAPTATION DYNAMICS
% =====================================================================

\paragraph{Why show the dual variable.}
The dual variable $\lambda_t$ is the central internal signal of the
BudgetPacer: it acts as a learned, time-varying ``price of cost''
that modulates the score of every arm at every step
(Eq.~\ref{eq:dual_update}).
When $\lambda_t$ is high, the router penalizes expensive arms
heavily; when $\lambda_t$ drops toward zero, cost is effectively
ignored and the router optimizes for quality alone.
Plotting $\lambda_t$ therefore exposes the \emph{mechanism} by
which the pacer adapts---not just the outcome (routing mix, cost)
but the control signal that produces it.
This is important because, unlike static cost penalties, $\lambda_t$
is never manually tuned: it emerges from the interaction between
the primal routing decisions and the dual budget feedback.
Observing its trajectory lets us verify that the adaptation is
smooth, timely, and regime-appropriate.

\paragraph{Lambda adaptation (Figure~\ref{fig:exp03_adaptation}a).}
The dual variable $\lambda_t$ responds directly to the cost EMA,
and all three budget regimes show the same qualitative Phase~2
response: $\lambda_t$ decays as post-drop costs fall below target.
The speed and depth of the decay depend on the margin between
post-drop costs and the budget target.
At \textbf{moderate} budget, $\lambda_t$ averages
$\bdParetoBanditModLambdaPhaseOne$ during Phase~1 then drops sharply
to ${\sim}\bdParetoBanditModLambdaPhaseTwo$ within Phase~2:
the post-drop portfolio cost (\bdParetoBanditModPhaseTwoCost/req)
is roughly half the target (\$6.6e-4), so the cost EMA falls well
below target and the dual update drives $\lambda_t$ toward
zero---the budget constraint becomes slack and the router reverts
to quality-only optimization.
At \textbf{loose} budget, the same mechanism acts even faster:
post-drop costs are only $\bdParetoBanditLoosePhaseTwoRatio$ the
target, and $\lambda_t$ drops from
$\bdParetoBanditLooseLambdaPhaseOne$ to near-zero almost
immediately.
At \textbf{tight} budget, $\lambda_t$ decays from
$\bdParetoBanditTightLambdaPhaseOne$ in Phase~1 to
${\sim}\bdParetoBanditTightLambdaPhaseTwo$ in Phase~2---slower
than moderate or loose because the margin is smaller.  Post-drop
costs (\bdParetoBanditTightPhaseTwoCost/req) are
$\bdParetoBanditTightPhaseTwoRatio$ the tight target (\$3.0e-4),
so the cost EMA settles just below target and $\lambda_t$
stabilizes at a small positive value rather than reaching zero.
This EMA-mediated response requires no external change-detection
signal: the dual update (Eq.~\ref{eq:dual_update}) naturally tracks
the cost regime.

\paragraph{Phase~3 lambda recovery.}
When Phase~3 restores original pricing, $\lambda_t$ rises back
toward Phase~1 levels, demonstrating full round-trip adaptation.
At \textbf{tight} budget, $\lambda_t$ climbs from
${\sim}\bdParetoBanditTightLambdaPhaseTwo$ to
${\sim}\bdParetoBanditTightLambdaPhaseThree$;
at \textbf{moderate}, from
${\sim}\bdParetoBanditModLambdaPhaseTwo$ to
${\sim}\bdParetoBanditModLambdaPhaseThree$;
at \textbf{loose}, from near-zero to
${\sim}\bdParetoBanditLooseLambdaPhaseThree$.
The recovery demonstrates that the primal-dual mechanism adapts
in both directions: $\lambda_t$ decreases when costs drop below
target and increases when costs rise above it.
Phase~3 $\lambda_t$ does not exactly match Phase~1 levels: at tight
budget, $\lambda_t$ overshoots ($0.51$ vs.\ $0.41$ in Phase~1),
while at loose budget it closely tracks ($0.10$ vs.\ $0.08$).
This asymmetry reflects the ramp-up transient---the dual update must
raise $\lambda_t$ from near-zero through the slack region---and is
consistent with the finite phase length (608 steps).

\paragraph{Gemini adoption (Figure~\ref{fig:exp03_adaptation}b).}
All three budget regimes show substantial Gemini adoption in Phase~2.
At \textbf{tight} budget, Gemini selection surges from
${\sim}\bdParetoBanditTightGeminiPhaseOne\%$ to
${\sim}\bdParetoBanditTightGeminiPhaseTwo\%$---the highest of any
constrained condition.
Two mechanisms contribute: the residual Phase~1 $\lambda_t$
(${\sim}\bdParetoBanditTightLambdaPhaseOne$) persists into early
Phase~2, penalizing Mistral's non-zero normalized cost and driving
initial Gemini concentration; as $\lambda_t$ decays to
${\sim}\bdParetoBanditTightLambdaPhaseTwo$, the router settles into
a Gemini-heavy but mixed policy, since Gemini is simultaneously
the cheapest and highest-quality arm post-drop.
At \textbf{moderate} budget, Gemini selection rises from
${\sim}\bdParetoBanditModGeminiPhaseOne\%$ to
${\sim}\bdParetoBanditModGeminiPhaseTwo\%$, approaching the
unconstrained baseline (${\sim}59\%$ Gemini in Phase~2).
At \textbf{loose} budget, Gemini selection grows from
${\sim}\bdParetoBanditLooseGeminiPhaseOne\%$ to
${\sim}\bdParetoBanditLooseGeminiPhaseTwo\%$, also converging
toward the unconstrained level.
The tight condition retains higher Gemini adoption than moderate or
loose because its residual $\lambda_t$
(${\sim}\bdParetoBanditTightLambdaPhaseTwo$) still provides mild
cost pressure favoring the cheapest arm, whereas at moderate and
loose $\lambda_t$ reaches near-zero and the router quality-optimizes
with a mixed Gemini/Mistral policy.

In Phase~3, Gemini adoption drops back to near Phase~1 levels:
tight $\bdParetoBanditTightGeminiPhaseThree\%$, moderate
$\bdParetoBanditModGeminiPhaseThree\%$, loose
$\bdParetoBanditLooseGeminiPhaseThree\%$.
The symmetric reversal confirms that the routing mix is driven by
the dual variable rather than by stale arm statistics.

\paragraph{Quality cost decomposition.}
The unconstrained baseline shares ParetoBandit's $\gamma{=}0.995$ but
removes the budget pacer, isolating the pacer's quality impact.
In Phase~2, the unconstrained baseline achieves reward $0.928$,
while ParetoBandit at moderate and loose budgets reaches
$0.927$---a gap of ${\leq}0.1\%$.
The pacer thus imposes \emph{negligible quality cost} when the
constraint is slack; at tight budget the cost is still only
${\sim}0.1\%$ ($\bdParetoBanditTightRewardPhaseTwo$ vs.\ $0.928$),
attributable to the residual $\lambda_t$ that slightly favors
cheaper arms over the quality-optimal mix.
%
The $\epsilon$-constraint procedure
(Section~\ref{sec:hparam_sweep}) selected $\gamma{=}0.995$ because
it trades a marginal stationary-quality loss for substantially
faster adaptation, and the present experiment validates that
tradeoff end-to-end: ParetoBandit's Phase~2 reward lift
($\Delta{=}+\bdParetoBanditLooseRewardLift$ to
$+\bdParetoBanditTightRewardLift$) confirms that the forgetting
factor enables the router to capitalise on the new cost landscape
rather than remaining anchored to stale Phase~1 statistics.

\paragraph{Budget compliance (Figure~\ref{fig:exp03_adaptation}c).}

\begin{center}
\small
\begin{tabular}{@{}lrrrr@{}}
\toprule
\textbf{Budget} & \textbf{Target} & \textbf{Phase~1} & \textbf{Phase~2} & \textbf{Phase~3} \\
\midrule
Tight    & \$3.00e-4 & \bdParetoBanditTightPhaseOneCost & \bdParetoBanditTightPhaseTwoCost & \bdParetoBanditTightPhaseThreeCost \\
Moderate & \$6.62e-4 & \bdParetoBanditModPhaseOneCost & \bdParetoBanditModPhaseTwoCost & \bdParetoBanditModPhaseThreeCost \\
Loose    & \$1.87e-3 & \bdParetoBanditLoosePhaseOneCost & \bdParetoBanditLoosePhaseTwoCost & \bdParetoBanditLoosePhaseThreeCost \\
\midrule
Unconstr.\ & ---     & \$9.5e-3 & \$3.2e-4 & \$9.8e-3 \\
\bottomrule
\end{tabular}
\end{center}

\noindent
ParetoBandit achieves Phase~1 costs within $3\%$ of the target at all
budget levels.  After the price drop, costs decline below target at
all three levels (the pacer constrains from above, never inflates):
tight costs settle at $\bdParetoBanditTightPhaseTwoRatio$ target,
moderate at $\bdParetoBanditModPhaseTwoRatio$, and loose at
$\bdParetoBanditLoosePhaseTwoRatio$.
In Phase~3, when original pricing is restored, ParetoBandit recovers
budget compliance to $1.00\times$--$1.02\times$ of the target across
all three regimes.
The unconstrained baseline spends ${\sim}5$--$33\times$ more than
pacer targets in Phase~1 (\$9.5e-3/req), confirming that the pacer
provides genuine cost control.

\paragraph{Reward lift after price drop.}
We test whether the pacer exploits the price drop by measuring
$\Delta_\mathrm{reward} = \bar{r}_\mathrm{Phase\,2} -
\bar{r}_\mathrm{Phase\,1}$ within ParetoBandit (a paired
comparison; each budget level serves as its own control).

\begin{center}
\small
\begin{tabular}{@{}lccc@{}}
\toprule
\textbf{Budget} & \textbf{P1 Rwd} & \textbf{P2 Rwd} & $\boldsymbol{\Delta}$ \\
\midrule
Tight    & \bdParetoBanditTightRewardPhaseOne & \bdParetoBanditTightRewardPhaseTwo & $+\bdParetoBanditTightRewardLift$ \\
Moderate & \bdParetoBanditModRewardPhaseOne & \bdParetoBanditModRewardPhaseTwo & $+\bdParetoBanditModRewardLift$ \\
Loose    & \bdParetoBanditLooseRewardPhaseOne & \bdParetoBanditLooseRewardPhaseTwo & $+\bdParetoBanditLooseRewardLift$ \\
\bottomrule
\end{tabular}
\end{center}

\noindent
All budget levels show highly significant reward improvement
($p < 10^{-3}$ for all; one-sample $t$-test on per-seed deltas,
19~d.f.).  The \textbf{tight} target exhibits the largest lift
($\Delta{=}+\bdParetoBanditTightRewardLift$), reflecting the shift
from a Llama/Mistral mix to
$\bdParetoBanditTightGeminiPhaseTwo\%$ Gemini routing once the
price drop makes Gemini the cheapest arm in the portfolio.
The moderate lift
($\Delta{=}+\bdParetoBanditModRewardLift$) is similar in magnitude;
the loose lift ($\Delta{=}+\bdParetoBanditLooseRewardLift$) is
smaller because the router was already partially selecting Gemini
in Phase~1.

\paragraph{Tight budgets: full round-trip adaptation.}
The tight target ($B{=}\$3.0{\times}10^{-4}$/req) is the most
instructive case because it is the only regime where the
constraint transitions from \emph{binding} to \emph{slack} and
back to \emph{binding} across the three phases.
In Phase~1, $\lambda_t$ averages
$\bdParetoBanditTightLambdaPhaseOne$, enforcing a Llama/Mistral
mix with negligible Gemini selection
($\bdParetoBanditTightGeminiPhaseOne\%$).
After the price drop, post-drop costs
(\bdParetoBanditTightPhaseTwoCost/req) fall below the target,
so $\lambda_t$ decays to
${\sim}\bdParetoBanditTightLambdaPhaseTwo$.
Despite the constraint becoming slack, the tight condition
still achieves the highest Gemini adoption
($\bdParetoBanditTightGeminiPhaseTwo\%$) of any constrained
condition, because the residual $\lambda_t$ provides mild cost
pressure that favors the now-cheapest arm (Gemini,
$\tilde{c}{\approx}0$) over the relatively expensive Mistral
($\tilde{c}{\approx}0.49$).
In Phase~3, when pricing is restored, $\lambda_t$ rises to
${\sim}\bdParetoBanditTightLambdaPhaseThree$, Gemini adoption
drops back to $\bdParetoBanditTightGeminiPhaseThree\%$, and the
budget recovers to $\bdParetoBanditTightPhaseThreeRatio$ of the
target.
This full round-trip---from aggressively binding Phase~1 to
mildly slack Phase~2 and back to binding Phase~3---demonstrates
that the primal-dual mechanism correctly tracks the constraint
status in both directions, without operator intervention.

\paragraph{Budget compliance across conditions
(Table~\ref{tab:budget_compliance}).}
Table~\ref{tab:budget_compliance} summarises cost/target ratios
for all conditions across all three phases.
ParetoBandit is the only condition that reliably tracks the budget
target in Phase~1 ($\bdParetoBanditLoosePhaseOneRatio$--$\bdParetoBanditTightPhaseOneRatio$).

After the Phase~2 price drop, the constraint becomes
\emph{non-binding} at all three budget levels: post-drop costs
fall below target for every condition, confirming that the
underspend is a property of the new cost landscape, not an
algorithmic failure.
At tight budget, the margin is small (ParetoBandit at
$\bdParetoBanditTightPhaseTwoRatio$ target), and $\lambda_t$
settles at a low but positive value
(${\sim}\bdParetoBanditTightLambdaPhaseTwo$).
At moderate and loose budgets, $\lambda_t$ decays to near-zero,
reverting to quality-only optimisation.

In Phase~3, when original pricing is restored, ParetoBandit
recovers to $\bdParetoBanditLoosePhaseThreeRatio$--$\bdParetoBanditModPhaseThreeRatio$
of the target---the only condition to achieve compliance in all
three phases.  By contrast, the Recalibrated Bandit (which
re-tunes its penalty at the Phase~3 boundary) reaches only
$\bdRecalTightPhaseThreeRatio$--$\bdRecalModPhaseThreeRatio$ at
tight and moderate budgets, while the Naive Bandit's Phase~3
compliance is erratic ($\bdNaiveModPhaseThreeRatio$--$\bdNaiveLoosePhaseThreeRatio$).

\begin{table}[t]
\centering
\caption{Budget compliance under cost drift (Experiment~3,
20~seeds, three phases).  Each cell shows realised average cost
as a multiple of the budget target ($1.00\times$ = perfect).
\textbf{Bold} marks values within $5\%$ of $1.00\times$.
$\dagger$~Phase~2 constraint non-binding: the price drop reduces
costs below target, regardless of algorithm.
ParetoBandit is the only condition that reliably meets the target in
Phase~1 ($\bdParetoBanditLoosePhaseOneRatio$--$\bdParetoBanditTightPhaseOneRatio$) and recovers compliance
in Phase~3 after pricing is restored.}
\label{tab:budget_compliance}
\small
\begin{tabular}{@{}llccc@{}}
\toprule
\textbf{Budget} & \textbf{Condition}
  & \textbf{Phase~1} & \textbf{Phase~2} & \textbf{Phase~3} \\
\midrule
\multirow{4}{*}{Tight ($\$3.0{\times}10^{-4}$)}
  & Fixed Policy        & $\bdFixedTightPhaseOneRatio$ & $\bdFixedTightPhaseTwoRatio^{\dagger}$ & $\bdFixedTightPhaseThreeRatio$ \\
  & Naive Bandit        & $\bdNaiveTightPhaseOneRatio$ & $\bdNaiveTightPhaseTwoRatio^{\dagger}$ & $\bdNaiveTightPhaseThreeRatio$ \\
  & Recalibrated        & $\bdRecalTightPhaseOneRatio$ & $\bdRecalTightPhaseTwoRatio$ & $\bdRecalTightPhaseThreeRatio$ \\
  & \textbf{ParetoBandit}  & $\mathbf{\bdParetoBanditTightPhaseOneRatio}$ & $\mathbf{\bdParetoBanditTightPhaseTwoRatio}$ & $\mathbf{\bdParetoBanditTightPhaseThreeRatio}$ \\[3pt]
\multirow{4}{*}{Moderate ($\$6.6{\times}10^{-4}$)}
  & Fixed Policy        & $\bdFixedModPhaseOneRatio$ & $\bdFixedModPhaseTwoRatio^{\dagger}$ & $\bdFixedModPhaseThreeRatio$ \\
  & Naive Bandit        & $\bdNaiveModPhaseOneRatio$ & $\bdNaiveModPhaseTwoRatio^{\dagger}$ & $\bdNaiveModPhaseThreeRatio$ \\
  & Recalibrated        & $\bdRecalModPhaseOneRatio$ & $\bdRecalModPhaseTwoRatio^{\dagger}$ & $\bdRecalModPhaseThreeRatio$ \\
  & \textbf{ParetoBandit}  & $\mathbf{\bdParetoBanditModPhaseOneRatio}$ & $\bdParetoBanditModPhaseTwoRatio^{\dagger}$ & $\mathbf{\bdParetoBanditModPhaseThreeRatio}$ \\[3pt]
\multirow{4}{*}{Loose ($\$1.9{\times}10^{-3}$)}
  & Fixed Policy        & $\bdFixedLoosePhaseOneRatio$ & $\bdFixedLoosePhaseTwoRatio^{\dagger}$ & $\bdFixedLoosePhaseThreeRatio$ \\
  & Naive Bandit        & $\bdNaiveLoosePhaseOneRatio$ & $\bdNaiveLoosePhaseTwoRatio^{\dagger}$ & $\bdNaiveLoosePhaseThreeRatio$ \\
  & Recalibrated        & $\bdRecalLoosePhaseOneRatio$ & $\bdRecalLoosePhaseTwoRatio^{\dagger}$ & $\bdRecalLoosePhaseThreeRatio$ \\
  & \textbf{ParetoBandit}  & $\mathbf{\bdParetoBanditLoosePhaseOneRatio}$ & $\bdParetoBanditLoosePhaseTwoRatio^{\dagger}$ & $\mathbf{\bdParetoBanditLoosePhaseThreeRatio}$ \\
\bottomrule
\end{tabular}
\end{table}

\paragraph{Synthesis.}
This experiment demonstrates that
\textbf{static cost penalties---even with periodic offline
recalibration---cannot substitute for principled budget pacing}:
\begin{itemize}[nosep,leftmargin=*]
  \item The \textbf{Fixed Policy} under-spends the budget by
    ${\sim}2\times$ in Phase~1 because it cannot refine its reward
    estimates---it sacrifices quality unnecessarily.
  \item The \textbf{Naive Bandit} is inconsistent: it overshoots
    at tight ($\bdNaiveTightPhaseOneRatio$) and moderate
    ($\bdNaiveModPhaseOneRatio$), and over-spends by
    $\bdNaiveLoosePhaseOneRatio$ at loose budget, because no
    feedback mechanism corrects the static cost penalty.
  \item The \textbf{Recalibrated Bandit} demonstrates that
    \emph{knowing} the optimal penalty at each phase boundary
    (via an offline sweep) is not enough: it inherits the Naive
    Bandit's Phase~1 non-compliance, and its stepwise recalibration
    cannot track intra-phase spending dynamics.
    In Phase~3, recalibration produces
    $\bdRecalTightPhaseThreeRatio$--$\bdRecalModPhaseThreeRatio$
    compliance at tight and moderate budgets---better than the Naive
    Bandit but far from target.
  \item \textbf{ParetoBandit} is the only condition that achieves
    precise Phase~1 budget compliance
    ($\bdParetoBanditLoosePhaseOneRatio$--$\bdParetoBanditTightPhaseOneRatio$ target)
    while automatically exploiting the Gemini price
    drop to improve quality
    ($\Delta_\mathrm{reward}{=}+\bdParetoBanditLooseRewardLift$ to
    $+\bdParetoBanditTightRewardLift$).
  \item \textbf{Phase~3 round-trip recovery:} when original pricing
    is restored, ParetoBandit recovers to
    $\bdParetoBanditLoosePhaseThreeRatio$--$\bdParetoBanditTightPhaseThreeRatio$
    of the budget target across all three regimes.  No other
    condition achieves compliance in all three phases---this
    demonstrates that the primal-dual mechanism adapts symmetrically
    to both cost decreases and cost increases.
  \item All three budget regimes demonstrate consistent
    \textbf{dual-variable dynamics} that follow from the Lagrangian
    structure: when post-drop costs fall below target, $\lambda_t$
    decays and the router shifts toward quality-only optimization;
    when Phase~3 restores costs above target, $\lambda_t$ rises
    and the router re-engages cost control.
    The decay speed depends on the margin---moderate and loose
    reach near-zero within Phase~2, while tight stabilizes at a
    small positive value
    (${\sim}\bdParetoBanditTightLambdaPhaseTwo$) because the margin
    is narrow.
    This is the correct Lagrangian behavior---the pacer
    automatically identifies the degree of constraint slack
    and adjusts accordingly.
\end{itemize}
